\documentclass[12pt]{article}
\usepackage[utf8]{inputenc}
\usepackage{amsmath}
\usepackage{geometry}
\usepackage{fancyhdr}
\geometry{letterpaper, margin=1in}
\pagestyle{fancy}
\lhead{EMCH 290 Extra Exercise \#1 Solutions}
\rhead{Page \thepage}
\title{\vspace{-0.75in}EMCH 290 Extra Exercise \#1 Solutions}
\date{\vspace{-0.3in}}

\begin{document}

\maketitle

\vspace{-0.5in}

\noindent Name: \underline{\hspace{3in}} \quad Date: \underline{\hspace{1.5in}}

\vspace{0.2in}

\subsection*{Solution to Question 1 [25 points]}
You are allowed to use internet conversion tables, given that sources are referenced correctly.
\begin{enumerate}
    \item Convert 3000 kJ to cal. 
    \[ 3000 \, \text{kJ} \times \left(\frac{1000 \, \text{cal}}{4.184 \, \text{kJ}}\right) = 717056.39 \, \text{cal} \]
    \item Convert 300 m/s to mph.
    \[ 300 \, \text{m/s} \times \left(\frac{3600 \, \text{s}}{1 \, \text{hr}}\right) \times \left(\frac{1 \, \text{mile}}{1609.34 \, \text{m}}\right) = 671.08 \, \text{mph} \]
\end{enumerate}
\newpage

\subsection*{Solution to Question 2 [25 points]}
\begin{enumerate}
    \item Convert 500 ft-lbf/min to W.
    \[ 500 \, \text{ft-lbf/min} \times \left(\frac{1 \, \text{W}}{0.7376 \, \text{ft-lbf/s}}\right) \times \left(\frac{1 \, \text{min}}{60 \, \text{s}}\right) = 11.29 \, \text{W} \]
    \item Convert 20000 ft/s² to m/s².
    \[ 20000 \, \text{ft/s}^2 \times \left(\frac{0.3048 \, \text{m}}{1 \, \text{ft}}\right) = 6096 \, \text{m/s}^2 \]
\end{enumerate}
\newpage

\subsection*{Solution to Question 3 [25 points]}
\[
\text{Weight on Mars} = \text{Weight on Earth} \times \left(\frac{g_{\text{Mars}}}{g_{\text{Earth}}}\right) 
= 2000 \, \text{N} \times \left(\frac{3.72076 \, \text{m/s}^2}{9.81 \, \text{m/s}^2}\right) = 759.8 \, \text{N}
\]
\newpage

\subsection*{Solution to Question 4 [25 points]}

First, let's write the formula for the total pressure at the bottom of the tank:
\[
P_{\text{bottom}} = P_{\text{water}} + P_{\text{oil}} + P_{\text{gasoline}}
\]
Next, we substitute in the formula for pressure \( P = \rho \times g \times h \) for each liquid layer:
\[
P_{\text{bottom}} = \rho_{\text{water}} \times g \times h_{\text{water}} + \rho_{\text{oil}} \times g \times h_{\text{oil}} + \rho_{\text{gasoline}} \times g \times h_{\text{gasoline}}
\]
Now we plug in the known values for each layer:
\begin{align*}
&= 1000 \, \text{kg/m}^3 \times 9.81 \, \text{m/s}^2 \times 2 \, \text{m} \\
&+ 890 \, \text{kg/m}^3 \times 9.81 \, \text{m/s}^2 \times 2 \, \text{m} \\
&+ 750 \, \text{kg/m}^3 \times 9.81 \, \text{m/s}^2 \times 2 \, \text{m}
\end{align*}
Calculating the pressure contributions from each layer:
\[
= 19620 \, \text{Pa} + 17442 \, \text{Pa} + 14705 \, \text{Pa}
\]
Finally, summing these up gives:
\[
= 51767 \, \text{Pa} = 51.8 \, \text{kPa}
\]


\end{document}
